Gliedern Sie die relevante Legislative von Deutschland und der EU

Förmliche Gesetze bzw. Gesetze im formellen Sinn werden vom parlamentarischen Gesetzgeber in dem in der Verfassung dafür
vorgesehenen Verfahren (Artikel 76 bis 82 Grundgesetz) beschlossen.
Davon zu unterscheiden sind Verordnungen, welche
Generell- verbindliche Rechtsnorm, welche durch Regierungs- oder Verwaltungsorgan erlassen wird.
Es bedarf einer Ermächtigungsgrundlage im Gesetz.
In der EU ist eine Verordnung ein Rechtsakt, der nach Verabschiedung durch die Mitgliedsstaaten unmittelbare Geltung hat –
sie muss nicht wie eine Richtlinie durch die nationalen Parlamente durch innerstaatliche Gesetze umgesetzt werden.

Das deutsche Recht untergliedert sich in drei Teilbereiche: Zivilrecht, Öffentli-
ches Recht und Strafrecht.

Das Zivilrecht (auch Privatrecht oder Bürgerliches Recht) regelt die Rechtsbezie-
hungen zwischen rechtlich gleichgeordneten Rechtssubjekten (natürliche und ju-
ristische Personen) untereinander. Zum Zivilrecht zählen beispielsweise das Versi-
cherungsvertragsgesetz (VVG), das Pflichtversicherungsgesetz (PflVG), das Bür-
gerliche Gesetzbuch (BGB), die VVG-Informationspflichtenverordnung (VVG-
InfoV), das Handelsgesetzbuch (HGB), das Aktiengesetz (AktG).

Das Öffentliche Recht regelt die Rechtsbeziehungen des Einzelnen (natürliche
und juristische Personen) zum Staat (Träger öffentlicher Gewalt) sowie die Rechts-
beziehungen der Verwaltungsträger untereinander. Zum Öffentlichen Recht zählen
beispielsweise das Versicherungsaufsichtsgesetz (VAG), die Aktuarverordnung
(AktuarV), die Anlageverordnung (AnlV), die Deckungsrückstellungsverordnung
(DeckRV), die Kapitalausstattungs-Verordnung (KapAusstV),

Das Strafrecht umfasst die Gesamtheit der Rechtsnormen, in denen die Voraus-
setzungen für eine Straftat und ihre Rechtsfolgen festgelegt sind. Das Strafrecht
ist insbesondere im Strafgesetzbuch (StGB) geregelt.

Primäres Gemeinschaftsrecht

Das Primärrecht ist mit dem Verfassungsrecht auf der Ebene der
Nationalstaaten vergleichbar. In den EU-Verträgen sind die genauen
Entscheidungsabläufe zur Setzung europäischen Rechts, das
institutionelle Gefüge und die Kompetenzen der Union geregelt.
Beispiel: Gründungs-, Revisions- und Beitrittsverträge.

Sekundäres Gemeinschaftsrecht
- Verordnungen (vergleichbar mit einem Gesetz haben unmittelbare Geltung in jedem Mitgliedstaat, z.B. Solvency II)
- Richtlinien (sind in nationales Recht noch umzusetzen, EU-Vermittlerrichtlinie, DSGVO)
- Empfehlungen und Stellungnahmen (nicht bindend)

Skizieren Sie die allgemeinen Bestandteile für das Zustandekommen eines Vertrages

Ein Vertrag ist ein zwei- oder mehrseitiges Rechtsgeschäft. Ein Vertrag kommt
durch (mindestens) zwei übereinstimmende Willenserklärungen zustande, An-
gebot (§ 145 BGB) und Annahme (§ 147 BGB).
Das Angebot ist eine empfangsbedürftige Willenserklärung, durch die einem an-
deren ein Vertragsschluss in der Weise angetragen wird, dass das Zustandekom-
men des Vertrages nur von dessen Einverständnis abhängt. Das Angebot muss so
bestimmt sein, dass die Annahme durch eine bloße Zustimmung durch einen an-
deren erfolgen kann, d.h. das Angebot muss alle wesentlichen Bestandteile des
Vertrages enthalten. Der Antragende ist für eine gewisse Zeit an seinen Antrag
gebunden (§ 145 BGB).
Die Annahme ist eine grundsätzlich empfangsbedürftige Willenserklärung, durch
die der Antragsempfänger dem Antragenden sein vorbehaltloses Einverständnis
mit dem angebotenen Vertragsschluss zu verstehen gibt.
§ 5 VVG enthält gegenüber dem BGB eine Spezialregelung (lex specialis) für den
Fall von Abweichungen zwischen dem Antrag oder bereits getroffenen Vereinba-
rungen und der Police (Rechtsfolge: Fiktion der Genehmigung des abweichenden
Versicherungsscheins, wenn kein Widerspruch) und weicht insoweit von der allge-
meinen Regelung in § 150 Abs. 2 BGB (Ablehnung verbunden mit einem neuen
Angebot) ab.

Was sind klassische Bestandteile des Versicherungsvertrages?

Der Versicherungsvertrag kommt ebenfalls durch Angebot und Annahme zu-
stande. Es gelten jedoch einige Sonderregeln des VVG als lex specialis bzw. er-
gänzend zum BGB. Dies sind auf Seiten des Versicherers (VR) die Beratungs-,
Informations- und Dokumentationspflichten und auf Seiten des Versicherungsneh-
mers (VN) die vorvertragliche Anzeigepflicht (Einzelheiten unter 2.3 und 2.4).
Der konkrete Inhalt eines Versicherungsvertrags folgt aus den (wirksamen) Ab-
reden der Vertragsparteien. Grundelemente eines Versicherungsvertrags sind re-
gelmäßig:
- Vertragsparteien (ggf. versicherte Personen oder Bezugsberechtigte)
- Leistungsversprechen (z. B. Definition des Versicherungsfalls)
- Versicherungssumme (ggf. Selbstbeteiligung)
- Prämienhöhe und -zahlweise
- Obliegenheiten (Verhaltenspflichten, die zum Erlöschen des Versicherungs-
schutzes führen können; z. B. vorvertragliche Anzeigepflicht)
- Risikoausschlüsse oder -zuschläge
- Vertragsbeginn und Vertragslaufzeit


Nennen Sie wesentlichen Rechte und Pflichten des Versicherers (VR)

Vor dem Vertragsschluss besteht für den VR eine Beratungspflicht (§ 6 VVG). 
Auf Grundlage der verbundenen Befragung schuldet er dann eine Beratung im 
Verhältnis zur Komplexität der Situation (anlassbezogene Beratung), welche dokumentiert werden muss.
Rechtsfolge einer Beratungspflichtverletzung ist insbesondere eine Verpflichtung 
zum Schadensersatz (§ 6 Abs. 5 VVG).

Der VR dem VN rechtzeitig vor Abgabe der Vertragserklärung
die Versicherungsbestimmungen inkl. AVB und die Informationen nach der VVG-
Informationspflichtenverordnung zur Verfügung zu stellen (Informationspflicht,
§ 7 VVG). Rechtsfolge fehlender Information ist insbesondere, dass die Wider-
rufsfrist nicht zu laufen beginnt (§ 8 Abs. 2 VVG).

Nach Vertragsschluss ist vertragstypische Verpflichtung des VR, ein be-
stimmtes Risiko des VN oder eines Dritten abzusichern (§ 1 Satz 1 VVG). Die
Absicherung geschieht (zunächst) durch das Versprechen einer Leistung bei Eintritt
des vereinbarten Versicherungsfalls und (sodann) durch die Leistung bei Eintritt
des Versicherungsfalls.

Der VR hat ein allgemeines Recht zur Prämienanpassung (§ 40 VVG). Vo-
raussetzung ist eine Prämienanpassungsklausel in den AVB. Erhöht der VR die Prä-
mie, steht dem VN ein Sonderkündigungsrecht zu. In der Lebens- und der Kran-
kenversicherung gelten Sonderregelungen zur Prämienanpassung (§ 163 Abs. 1
bzw. § 203 Abs. 2 VVG).

Nennen Sie wesentliche Rechte und Pflichten des Versicherungsnehmers (VN)

Der VN hat bis zur Abgabe seiner Vertragserklärung die ihm bekannten Ge-
fahrumstände, die für den Entschluss des VR, den Vertrag mit dem vereinbarten
Inhalt zu schließen, erheblich sind und nach denen der VR in Textform gefragt hat,
dem VR anzuzeigen (vorvertragliche Anzeigepflicht, § 19 Abs. 1 VVG).
Verletzt der VN seine vorvertragliche Anzeigepflicht, ist der VR grundsätzlich zum Rücktritt vom Vertrag
berechtigt (§ 19 Abs. 2 VVG). Unter bestimmten Voraussetzungen tritt an die Stelle
des Rücktritts- ein Kündigungs- oder ein Vertragsanpassungsrecht (§ 19 Abs. 3 und 4 VVG).


Grundsätzlich besteht bei jedem Versicherungsvertrag ein Widerrufsrecht
(§ 8 VVG). Die zweiwöchige Widerrufsfrist beginnt mit Zugang von Versicherungsschein, 
AVB und Informationen nach der VVG-Informationspflichtenverordnung sowie 
Belehrung zu laufen (§ 8 Abs. 2 VVG).

Nach Vertragsschluss ist vertragstypische Verpflichtung des VN die Prä-
mienzahlung (§ 1 Satz 2 VVG).
Hinsichtlich der Rechtsfolge bei Verletzung der Prämienzahlungspflicht ist zwischen
Erst- (§ 37 VVG) und Folgeprämie (§ 38 VVG) zu differenzieren

Für den VN bestehen gesetzliche (z.B. Gefahrstandsobliegenheit, siehe unter
2.4.5) und vertragliche Obliegenheiten. Die Rechtsfolgen der Verletzung einer
vertraglichen Obliegenheit sind in § 28 VVG geregelt. Die Vorschrift unterscheidet
bezüglich der Rechtsfolgen von Obliegenheitsverletzungen zwischen Obliegen-
heitsverletzungen vor und nach Eintritt des Versicherungsfalls. Geregelt sind das
Schicksal des Vertrags und die Leistungspflicht des VR (ggf. Leistungskürzung oder
Leistungsfreiheit)

Eine Gefahrerhöhung ist jede vom status quo bei Antragstellung abweichende erhebliche, nicht vorhersehbare Änderung, 
die die Wahrscheinlichkeit eines Schadenseintritts oder die Vergrößerung des Schadens erhöht.

Der VN darf eine Gefahrerhöhung nicht selbst herbeiführen (§ 23 Abs. 1 VVG) und
hat sowohl eine unbewusst selbst herbeigeführte (subjektive) als auch eine ohne
sein Zutun herbeigeführte (objektive) Gefahrerhöhung unverzüglich anzuzeigen (§
23 Abs. 2 und 3 VVG). Kumulative Voraussetzungen sind (1.) Erheblichkeit der
Gefahrerhöhung, (2.) Unvorhersehbarkeit der Gefahrerhöhung und (3.) eine nicht
nur vorübergehende Änderung.

Der VR kann bei einer Gefahrerhöhung abhängig vom Grad des Verschuldens (Vor-
satz, grobe Fahrlässigkeit, einfache Fahrlässigkeit) den Vertrag nach seiner Wahl
kündigen (§ 24 VVG) oder eine Prämienerhöhung bzw. einen Risikoausschluss vor-
nehmen (§ 25 VVG). Für die Frage nach der Leistungspflicht des VR (ggf. Leis-
tungskürzung oder Leistungsfreiheit) ist ebenfalls der Grad des Verschuldens ent-
scheidend (§ 26 VVG).


Differenzieren sie die Arten von Versicherungsvermittlern und deren Vorraussetzungen

Versicherungsvermittler ist der Oberbegriff für Versicherungsmakler und Versi-
cherungsvertreter. Versicherungsvermittler kann eine natürliche Person, eine Per-
sonengesellschaft oder eine juristische Person sein.

Versicherungsvertreter ist (§ 59 Abs. 2 VVG), wer von
einem Versicherer oder einem Versicherungsvertreter damit betraut ist, gewerbs-
mäßig Versicherungsverträge zu vermitteln oder abzuschließen. Der Versiche-
rungsvertreter steht im Lager des VR, ist Erfüllungsgehilfe des VR und besitzt ge-
mäß §§ 69 – 73 VVG gesetzliche Vertretungsmacht für den VR („Auge und Ohr“
des VR). Zu den Versicherungsvertretern im Sinne des Gesetzes zählen Ausschließ-
lichkeitsvertreter, Mehrfachvertreter, aber auch Angestellte des VR im Außen-
dienst.

Versicherungsmakler ist (§ 59 Abs. 3 Satz 1 VVG), wer
gewerbsmäßig für den Auftraggeber die Vermittlung oder den Abschluss von Versicherungsverträgen übernimmt, 
ohne von einem Versicherer oder von einem Versicherungsvertreter damit betraut zu sein. 
Der Versicherungsmakler steht im Lager des VN als dessen Sachwalter und Interessenwahrer. 
Er ist nicht nur zur einmaligen Beschaffung von Versicherungsschutz verpflichtet, 
sondern auch zur anschließenden Dauerbetreuung mit Anpassung und Veränderung sowie Verwaltung der
Verträge.

Wer gewerbsmäßig als Versicherungsmakler oder als Versicherungsvertreter den
Abschluss von Versicherungsverträgen vermitteln will, bedarf der Erlaubnis der
zuständigen Industrie- und Handelskammer (§ 34d Abs. 1 Satz 1 GewO).
Kumulative Voraussetzungen der Erlaubniserteilung sind (§ 34d Abs. 1 und 2
GewO):
(1.)Zuverlässigkeit
(2.)Geordnete Vermögensverhältnisse
(3.)Berufshaftpflichtversicherung
(4.)Sachkundeprüfung

Eine uneingeschränkte Erlaubnispflicht gilt für Nicht-Ausschließlichkeitsvertre-
ter und Makler (§ 34d Abs. 1 GewO). Von der Erlaubnispflicht befreit sind sog.
produktakzessorische Vermittler (§ 34d Abs. 3 GewO), z.B. Kfz-Haftpflicht- und
Kaskoversicherung beim Kfz-Kauf.

Erlaubnisfrei sind sog. gebundene Vermittler unter den folgenden zwei Vo-
raussetzungen(§ 34d Abs. 4 GewO):
(1.)Ausübung der Tätigkeit ausschließlich im Auftrag eines oder, wenn die Ver-
sicherungsprodukte nicht in Konkurrenz stehen, mehrerer im Inland zum
Geschäftsbetrieb befugten Versicherungsunternehmen
(2.)Uneingeschränkte Haftungsübernahme für die Vermittlertätigkeit durch das
oder die Versicherungsunternehmen


Erläutern Sie Details zur Schaden und Summenversicherung 

In der Schadensversicherung ersetzt der VR dem VN im Versicherungsfall den in
der Höhe entstandenen (Schaden-)Betrag (z.B. Kfz-Versicherung). 
In der Summenversicherung, wird die im Vertrag vereinbarte Summe unabhängig
vom tatsächlich entstandenen Schaden geleistet (z.B. Risikolebensversicherung).

Für die Schadensversicherung regelt § 81 VVG die Herbeiführung des Ver-
sicherungsfalls. Führt der VN den VersFall vorsätzlich („Wissen und Wollen“,
keine Definition im BGB) herbei, ist der VR leistungsfrei.
Im Falle grob fahrlässiger Herbeiführung des VersFalls ist der VR zur Leistungskürzung in einem der Schwere
des Verschuldens des VN entsprechenden Verhältnis berechtigt.
Hinsichtlich der Beweislastverteilung hat der VN den Eintritt des VersFalls zu be-
weisen. Der VR hat die Herbeiführung des VersFalls, Kausalität und Verschulden
darzulegen und ggf. nachzuweisen.

Bei Eintritt des Versicherungsfalls hat der VN nach § 82 VVG für die Abwendung
und Minderung des Schadens zu sorgen. Der VN soll durch die sogen. Rettungsobliegenheit dazu angehalten werden,
und kann seine Rettungskosten (Aufwendungen) insoweit zu erstatten, als der VN sie
für geboten halten durfte. Grundsätzlich sind solche Kosten nur zulässig wenn sie während/nach oder unmittelbar vor Schadeneintritt entstehen.

Erleidet der VN einen Schaden, gegen den er versichert ist, durch die schädi-
gende Handlung eines Dritten, hat er die Wahlmöglichkeit, seinen Schaden entwe-
der vom Schädiger ersetzt zu verlangen oder seinen VR in Anspruch zu nehmen.
Die Vorschrift des § 86 VVG regelt den Übergang des Ersatzanspruchs auf den
vom VN/Geschädigten in Anspruch genommenen VR.
Der VR darf seinen Regressanspruch erst verwirklichen, wenn der VN insgesamt
seinen Schaden ersetzt bekommen hat (§ 86 Abs. 1 Satz 2 VVG, sog. Quotenvor-
recht). Der VN hat die Obliegenheit zur Wahrung des Ersatzanspruchs und zur
Mitwirkung bei dessen Durchsetzung (§ 86 Abs. 2 VVG). Der VR hat keinen durch-
setzbaren Rückgriffsanspruch gegen die mit dem VN in häuslicher Gemeinschaft
lebenden Personen (§ 86 Abs. 3 VVG – Ausnahme: Vorsatz).


Wie unterteilt sich Aufsicht im Wirtschaftsrecht und was sind ihre Ziele für Versicherungen 

Die vertikale Aufsicht beobachtet die gesamte Wirtschaft laufend un-
ter einem bestimmten Aspekt (z.B. Kartellaufsicht). Wohingegen die branchenspe-
zifische Form der Aufsicht als horizontale Aufsicht bezeichnet wird (z.B. Banken-
, Wertpapier-, Versicherungsaufsicht).

Ziele des Versicherungsaufsichtsrechts sind:
- Ausreichende Wahrung der Belange der Versicherten (§§ 11, 294 VAG)
- Dauernde Erfüllbarkeit der Versicherungsverträge (§§ 11, 294 VAG)
- Schutz der kollektiven Verbraucherinteressen (§ 4 Abs. 1a FinDAG)
- Schutz der Funktionsfähigkeit des Versicherungswesens (<-> Stabilität Finanzsystems)


Nennen Sie die Subjekte der Versicherungsaufsicht und deren Instrumente

Zum Kreis der Beaufsichtigten zählen:
- Versicherungsunternehmen (§ 1 Abs. 1 Nr. 1, § 7 Nr. 33 VAG)
- Versicherungs-Holdinggesellschaften (§ 1 Abs. 1 Nr. 2, § 7 Nr. 31 VAG)
- Versicherungszweckgesellschaften (§ 1 Abs. 1 Nr. 3, § 168 VAG)
- Sicherungsfonds (§ 1 Abs. 1 Nr. 4, § 223 VAG)
- Pensionsfonds (§ 1 Abs. 1 Nr. 5, § 236 Abs. 1 VAG)
- Inhaber bedeutender Beteiligungen (§§ 16 ff. VAG)

Die Aufsichtsbehörde nimmt die ihr zugewiesenen Aufgaben „nur im öffentlichen
Interesse“ wahr (§ 294 Abs. 8 VAG). Das hat zur Folge, dass eine Amtshaftung
gegenüber den Versicherten ausgeschlossen ist.

Aufgabe der Aufsichtsbehörde ist die Überwachung des gesamten Geschäftsbe-
triebs des Versicherungsunternehmens (§ 294 Abs. 2 Satz 1 VAG).
Die BaFin hat zwei Arten von Befugnissen: Die präventive Aufsicht umfasst ins-
besondere die Erlaubnis zum Geschäftsbetrieb, die Genehmigung von Geschäfts-
planänderungen und Bestandsübertragungen. Zur repressiven Aufsicht zählen
die speziellen Befugnisse nach §§ 298 ff. VAG (z.B. Einsetzung eines Sonderbe-
auftragten, § 307 VAG; Verwarnung und Abberufung von Geschäftsleitern und Per-
sonen mit Schlüsselaufgaben, § 303 VAG).

Was ist bei der Spartentrennung in Versicherungsaufsichtsrechts zu beachten?

Gemäß § 8 Abs. 4 S. 2 VAG darf nicht jedes Versicherungsunternehmen jede Versicherungssparte
betreiben.
- Übersicht der erlaubten Sparten findet sich in Anlage 1 zum VAG: Die Spartentrennung ist Pflicht
für Versicherungen, Lebensversicherungen, Krankenversicherungen sowie Schaden- und
Unfallversicherungen in rechtlich eigenständigen Unternehmen zu betreiben.
- Die jeweilige Unternehmensform muss eine Aktiengesellschaft, eine VVaG oder eine
Körperschaft/Anstalt öffentlichen Rechts sein.
- Die Spartentrennung führt im Versicherungsbereich zur Konzernbildung.
§ 8 Abs. 4 VAG: Ein Rückversicherungsunternehmen wird nur zum Betrieb der Rückversicherung
zugelassen.
Bei Erstversicherungsunternehmen schließen die Erlaubnis zum Betrieb der Lebensversicherung im
Sinne der Anlage 1 Nummer 19 bis 24 und die Erlaubnis zum Betrieb anderer Versicherungssparten
einander aus; das Gleiche gilt für die Erlaubnis zum Betrieb der Krankenversicherung im Sinne des
§ 146 Absatz 1 und die Erlaubnis zum Betrieb anderer Versicherungssparten


Was sind die gesetzlichen Anforderungen an das Risikomanagement und die versicherungsmathematische Funktion?

Nach § 26 Abs. 1 Satz 1 VAG müssen Versicherungsunternehmen über ein wirksa-
mes Risikomanagementsystem verfügen, das die Informationsbedürfnisse der
Personen, die das Unternehmen tatsächlich leiten oder andere Schlüsselfunktionen
innehaben, durch eine angemessene interne Berichterstattung gebührend berück-
sichtigt.
Zu den zu entwickelnden Strategien zählt insbesondere eine auf die Steuerung des
Unternehmens abgestimmte Risikostrategie, die Art, Umfang und Komplexität des
betriebenen Geschäfts und der mit ihm verbundenen Risiken berücksichtigt (§ 26
Abs. 2 VAG). Das Risikomanagementsystem hat sämtliche Risiken des Unterneh-
mens zu erfassen, wobei das Gesetz einen Mindestkatalog an Bereichen nennt, zu
denen innerbetriebliche Leitlinien zu erstellen sind (§ 26 Abs. 5 VAG).

Die versicherungsmathematische Funktion ist nunmehr obligatorisch (§ 31
Abs. 1 Satz 1 VAG). Der versicherungsmathematischen Funktion können – 
vorbehaltlich - der Vereinbarkeit mit dem Grundsatz angemessener Funktionstrennung
und den gesetzlichen Zuständigkeitszuweisungen – neben den in § 31 Abs. 1,2 VAG 
genannten Aufgaben 
- Berechnung der versicherungstechnischen Rückstellungen; 
- Mitwirkung bei Umsetzung des Risikomanagements
weitere Aufgaben zugewiesen werden. 

Denkbar ist daher – und bei kleineren Versicherungsunternehmen sinnvoll –, 
die versicherungsmathematische Funktion durch den in 
- Lebensversicherung
- private Krankenversicherung 
- Unfallversicherungen mit Prämienrückgewähr 
gesetzlich zwingend zu bestellenden Verantwortlichen Aktuar (§§ 141, 156 VAG) wahrnehmen zu lassen.

Der Verantwortliche Aktuar ist gesetzlich verpflichtet, seine Aufgaben als Verantwortlicher
Aktuar unabhängig von seinen sonstigen Pflichten auszuführen und darf keine Interessenskonflikte (z.B. Überschussbeteiligung) haben. 
Im Rahmen der Ausführung dieser Aufgaben ist er nicht an Weisungen gebunden, auch nicht seitens der
Unternehmensführung oder seines sonst für ihn zuständigen Vorgesetzten.
